\documentclass[10pt,a4paper]{article}

\usepackage[top=0.38in,bottom=0.38in,left=0.46in,right=0.46in]{geometry}
\usepackage[french]{babel}
\usepackage{fontspec}
\usepackage{xcolor}
\usepackage{hyperref}
\usepackage{graphicx}
\usepackage{enumitem}
\usepackage{tikz}
\usepackage{microtype}
\usepackage{titlesec}

\definecolor{linkblue}{HTML}{0B57D0}
\hypersetup{
  colorlinks=true,
  urlcolor=linkblue,
  linkcolor=linkblue
}

\pagestyle{empty}
\flushbottom
\linespread{1.00}
\pretolerance=10000
\tolerance=4000
\emergencystretch=3em
\hyphenpenalty=10000
\exhyphenpenalty=10000
\setlength{\parindent}{0pt}
\setlength{\parskip}{0.14em plus 0.04em minus 0.02em}
\setlist[itemize]{label=\textbullet, leftmargin=1.08em, itemsep=0.12em, topsep=0.14em plus 0.04em minus 0.02em, parsep=0em, partopsep=0em, before=\raggedright, after=\vspace{0.04em plus 0.02em minus 0.01em}}
\titleformat{\section}{\large\bfseries}{}{0em}{}[\titlerule]
\titlespacing*{\section}{0pt}{0.68em plus 0.12em minus 0.03em}{0.22em plus 0.04em minus 0.01em}
\graphicspath{{./}{../images/}}
\defaultfontfeatures{Ligatures=TeX}
\setmainfont{TeX Gyre Heros}[Scale=0.95]

\newcommand{\cvheading}[4]{
  \textbf{#1} \hfill #2 \\[0.05em]
  \textit{#3} \hfill \textit{#4}
}
\newcommand{\cvheadinglink}[5]{
  \textbf{\href{#5}{#1}} \hfill #2 \\[0.05em]
  \textit{#3} \hfill \textit{#4}
}
\newcommand{\contactlink}[2]{\href{#1}{\textcolor{linkblue}{\textbf{#2}}}}
\newcommand{\roundphoto}[1]{%
  \begin{tikzpicture}
    \clip (0,0) circle (0.975cm);
    \node at (0,0) {\includegraphics[width=1.95cm,height=1.95cm]{#1}};
  \end{tikzpicture}%
}

\begin{document}

\noindent
\begin{minipage}[t]{0.80\textwidth}
\vspace*{-0.58cm}
{\LARGE \textbf{Walid Gourideche}}\\[0.30em]
\small
Cite Al Amane Imm A Apt 1, 35000 Taza, Maroc\\[0.12em]
+212 698 585 528 \textbar{} \contactlink{mailto:walidgourideche@gmail.com}{walidgourideche@gmail.com}\\[0.10em]
\contactlink{https://grydsh.dev}{Portfolio} \textbar{} \contactlink{https://www.linkedin.com/in/walidgourideche/}{LinkedIn} \textbar{} \contactlink{https://github.com/Gryddd}{GitHub}
\end{minipage}
\hfill
\begin{minipage}[t]{0.16\textwidth}
\raggedleft
\raisebox{-1.04cm}[0pt][0pt]{\roundphoto{cv-photo.jpg}}
\end{minipage}
\vspace{0.06em}

\section*{Formation}
\cvheading{Institut Spécialisé de Technologie Appliquée (ISTA)}{sept. 2024 -- juil. 2026}{Technicien spécialisé en réseaux et systèmes informatiques}{Taza, Maroc}
\begin{itemize}
  \item Participation à des séances animées par les étudiants sur Linux, la virtualisation et l'administration de Windows Server.
\end{itemize}

\cvheading{Clinique de Hann. Münden}{oct. 2023 -- févr. 2024}{Période d'orientation de six mois en formation infirmière}{Hann. Münden, Allemagne}
\begin{itemize}
  \item Approfondissement de l'allemand et découverte de la culture de travail allemande.
\end{itemize}

\cvheading{Faculté Polydisciplinaire de Taza}{sept. 2021 -- juin 2022}{Études fondamentales en mathématiques et informatique}{Taza, Maroc}
\begin{itemize}
  \item Application de la logique mathématique à la résolution de problèmes en algorithmique et en informatique.
\end{itemize}

\cvheading{Lycée Al Falah 2}{sept. 2020 -- juin 2021}{Baccalauréat, option physique}{Taza, Maroc}
\begin{itemize}
  \item Bases solides en mathématiques et en physique utiles pour l'informatique et les réseaux.
\end{itemize}

\section*{Expérience professionnelle}
\cvheading{Société Régionale Multiservices Fès-Meknès (SRM)}{avr. 2026 -- mai 2026}{Stagiaire en réseau informatique}{Taza, Maroc}
\begin{itemize}
  \item Stage confirmé axé sur l'exploitation réseau, le support et des tâches de sécurité de base dans une infrastructure réglementée.
\end{itemize}

\cvheading{IQRA Société}{jan. 2026 -- févr. 2026}{Stagiaire en réseau informatique}{Casablanca, Maroc}
\begin{itemize}
  \item Supervision et administration de systèmes et services réseau via SSH et RDP dans un environnement d'entreprise.
\end{itemize}

\cvheading{Click Service}{juin 2025 -- sept. 2025}{Stagiaire en support informatique}{Taza, Maroc}
\begin{itemize}
  \item Installations logicielles et systèmes avec 98\% de réussite au premier passage, amélioration de 20\% de la précision de l'inventaire, mise en place ou maintenance de plus de 10 réseaux locaux et réduction de 25\% des demandes de support après installation.
\end{itemize}

\cvheading{Ministère de la Santé}{oct. 2022 -- déc. 2022}{Stagiaire en administration système}{Taza, Maroc}
\begin{itemize}
  \item Participation aux déploiements Hyper-V, réduisant de 30\% le temps de mise en place des environnements de test, et traitement de plus de 50 demandes liées à Active Directory.
\end{itemize}

\cvheading{Résidence pour seniors Inc.}{mai 2022 -- juil. 2022}{Stagiaire en conseil informatique}{Taza, Maroc}
\begin{itemize}
  \item Installation et configuration de matériel et d'équipements réseau.
\end{itemize}

\section*{Projets informatiques}
\cvheadinglink{LeagueSkins}{En cours}{Maintenance d'un projet open source et analyse binaire}{À distance}{https://github.com/Alban1911/LeagueSkins}
\begin{itemize}
  \item Maintenance de la base de ressources pour une application open source utilisée par plus de 10\,000 utilisateurs quotidiens dans le monde, rétro-ingénierie de fichiers binaires et coordination des tests de régression avant publication.
\end{itemize}

\cvheadinglink{PortGuardian Enterprise}{2026}{Système de détection et de réponse aux menaces USB}{Projet cybersécurité}{https://github.com/Gryddd/PortGuardian}
\begin{itemize}
  \item Technologies : Python, PyQt6, WMI et Splunk Syslog ; mise en œuvre d'une détection locale des périphériques USB, de l'analyse des fichiers, de l'isolation automatique et de la transmission des alertes pour des postes Windows, validée avec un temps moyen de détection USB de 350 ms, plus de 95\% de détection, moins de 5\% de faux positifs sur 200 échantillons, 98\% de réussite d'éjection et 99,8\% de fiabilité d'envoi vers Splunk.
\end{itemize}

\section*{Compétences techniques}
\textbf{Réseaux :} Cisco IOS, routage et commutation, TCP/IP, sous-réseaux, VLAN, DNS/DHCP, Wireshark, GNS3 \\
\textbf{Administration système :} Windows Server, Active Directory, Linux (Ubuntu/Debian), Hyper-V, VMware \\
\textbf{Sécurité :} pfSense, Suricata (IPS/IDS), pfBlockerNG, Nmap, analyse des vulnérabilités, ajustement des règles

\section*{Langues}
Allemand (Goethe B2), anglais (C2), français (courant), arabe (langue maternelle)

\end{document}
